\documentclass[12pt,a4paper]{article}
\usepackage[utf8]{inputenc}
\usepackage[T1]{fontenc}
\usepackage[brazil]{babel}
\usepackage{indentfirst}
\usepackage[margin=2.5cm]{geometry}

\title{Resenha Cr\'itica: Artigo 6 ---\\
The Reactive Manifesto 2.0}
\author{Gustavo Pessoa Firmino Duarte}
\date{5 de Abril de 2026}

\begin{document}
\maketitle

\section{Introdu\c{c}\~ao}

Publicado em 16 de setembro de 2014, o \textit{Reactive Manifesto} em sua vers\~ao
2.0 apresenta um conjunto de princ\'ipios para o design de sistemas distribu\'idos
modernos. Elaborado por um grupo de especialistas liderado por Jonas Bon\'er, Dave
Farley, Roland Kuhn e Martin Thompson, o documento identifica quatro propriedades
fundamentais que sistemas \textit{reativos} devem apresentar: \textbf{Responsivos}
(\textit{Responsive}), \textbf{Resilientes} (\textit{Resilient}),
\textbf{El\'asticos} (\textit{Elastic}) e \textbf{Orientados a Mensagens}
(\textit{Message-Driven}). Escrito em um momento em que a demanda por sistemas de
larga escala, alta disponibilidade e baixa lat\^encia se intensificava, o manifesto
articula uma vis\~ao coerente de como arquiteturas distribu\'idas devem ser
constru\'idas para enfrentar esses desafios.

\section{As Quatro Propriedades Reativas}

A propriedade central \'e a \textbf{responsividade}: um sistema reativo responde
a tempo, mantendo uma qualidade de servi\c{c}\~ao consistente mesmo sob carga ou
falha. Esse comprometimento com lat\^encias bound\'arias n\~ao \'e apenas uma
quest\~ao de desempenho; \'e tamb\'em um meio de detectar problemas rapidamente,
pois timeouts s\~ao muito mais informativos do que respostas tardias e
imprevis\'iveis.

A \textbf{resili\^encia} diz respeito \`a capacidade de o sistema permanecer
responsivo diante de falhas. O manifesto prop\~oe que falhas sejam tratadas por
replica\c{c}\~ao, isolamento e delega\c{c}\~ao: cada componente executa em
isolamento para que suas falhas n\~ao se propaguem, e a recupera\c{c}\~ao \'e
delegada a um supervisor externo ao componente falho. Isso contrasta com o modelo
tradicional de tratamento sincrono de erros, onde uma exce\c{c}\~ao em um ponto
pode derrubar toda a pilha de chamadas.

A \textbf{elasticidade} garante responsividade sob varia\c{c}\~oes de carga. Um
sistema el\'astico escala horizontalmente, sem pontos centrais de contenção
(\textit{bottlenecks}), adicionando ou removendo recursos conforme a demanda.
Isso pressup\~oe que os componentes n\~ao compartilhem estado mut\'avel e que
possam ser instanciados e encerrados de forma transparente.

Por fim, ser \textbf{orientado a mensagens} \'e o meio que viabiliza as tr\^es
propriedades anteriores: a troca de mensagens ass\'incronas e n\~ao-bloqueantes
cria fronteiras claras entre componentes, permite a localiza\c{c}\~ao transparente
(o componente n\~ao precisa saber onde o destinat\'ario est\'a fisicamente),
poss\'ibilita \textit{back-pressure} (a capacidade de o receptor sinalizar que
n\~ao consegue acompanhar o ritmo do emissor) e suporta o isolamento de falhas
por meio de filas de mensagens que absorvem picos de demanda.

\section{Conex\~ao com a Pr\'atica Profissional}

Durante meu est\'agio de Sistemas na ArcelorMittal, tive contato com aplica\c{c}\~oes
enterprise de suporte internacional que precisavam funcionar continuamente em
m\'ultiplos fus\'os hor\'arios e que integravam sistemas legados de diferentes
pa\'ises. A discuss\~ao sobre resili\^encia e mensageria no manifesto ressoou
fortemente com problemas reais que observei: chamadas s\'incronas entre sistemas
que, quando um servi\c{c}o ficava indispon\'ivel, cascateavam falhas para os sistemas
dependentes. A introdu\c{c}\~ao de filas de mensagens (como as que o manifesto
recomenda) em partes do pipeline era justamente uma das solu\c{c}\~oes discutidas
pela equipe para isolar esses pontos de falha.

No meu MVP de e-commerce, trabalhei com um modelo essencialmente s\'incrono:
requisi\c{c}\~ao HTTP $\to$ servi\c{c}o $\to$ banco de dados $\to$ resposta.
Embora funcional para o escopo do projeto, reconheço que essa arquitetura seria
inadequada em cen\'arios de alta concorr\^encia. Um pagamento processado de forma
s\'incrona, por exemplo, poderia bloquear outras requisi\c{c}\~oes enquanto aguarda
a resposta de um gateway externo --- exatamente o problema que a orienta\c{c}\~ao
a mensagens busca eliminar.

\section{Conclus\~ao}

O \textit{Reactive Manifesto} oferece mais do que um conjunto de t\'ecnicas;
oferece um modelo mental coerente para sistemas distribu\'idos modernos. Suas
quatro propriedades --- responsividade, resili\^encia, elasticidade e orienta\c{c}\~ao
a mensagens --- est\~ao profundamente interligadas, e a aus\^encia de qualquer uma
compromete as demais. O manifesto ajudou a popularizar padr\~oes como atores
(\textit{Akka}), \textit{event streaming} (Kafka) e arquiteturas baseadas em eventos,
hoje presentes em praticamente toda aplica\c{c}\~ao de larga escala. Para quem
projeto sistemas distribu\'idos, o document tem o m\'erito de articular em termos
claros o que antes era conhecimento t\'acito de especialistas em sistemas de alto
desempenho.

\end{document}
